\subsection{Bounded-dimensional radial patches}
\label{sec:fixed-d-patches}
\label{sec:sharp-high-dimensional-edge}

This subsection proves the fixed-dimensional spatial exclusion needed for the
\(d\ge3\) support law in Theorem~\ref{thm:main-support-law}. The moving radial
shell identifies the candidate KKT patches, and the uniform estimates in
Subsection~\ref{sec:sharp-high-dimensional-estimates} exclude every other
score location.

Let \(Z_1,\ldots,Z_n\) be i.i.d.\ \(\cN(0,I_d)\), \(d\ge3\), and set
\[
    H_n(t)
    =
    \frac1n\sum_{i=1}^n
    \exp\{t\cdot Z_i-|t|^2/2\},
    \qquad t\in\R^d .
\]

\subsubsection{The moving radial shell}

Let
\[
    Y_i=\frac{|Z_i|^2}{2},\qquad U_i=\frac{Z_i}{|Z_i|}\in S^{d-1},
    \qquad k=\frac d2 .
\]
Since \(Y_i\) has the \(\Gamma(k,1)\) law,
\[
    b_{n,d}
    =
    \log n+(k-1)\log\log n-\log\Gamma(k),
\]
the fixed-dimensional specialization of the Gamma extreme-value theory in
Section~\ref{sec:mv-gamma-normalization} gives the marked radial limit.
\begin{lemma}
\label{lem:radial-gamma-extremes}
The extremal process
\[
    \Xi_n=\sum_{i=1}^n\delta_{(Y_i-b_{n,d},\,U_i)}
\]
converges to a Poisson point process on \(\R\times S^{d-1}\) with intensity
\(e^{-x}dx\,\sigma_{d-1}(du)\), where \(\sigma_{d-1}\) is normalized surface
measure.  In particular, \(M_n=\max_i(Y_i-b_{n,d})\) satisfies
\[
    \Pbb\{M_n\le x\}\to\exp\{-e^{-x}\}.
\]
\end{lemma}

\begin{proof}
Proposition~\ref{prop:gamma-max} and
Lemma~\ref{lem:fixed-d-normalization-bridge} give the unmarked radial limit.
Radius and direction are independent, and \(U_i\) has law
\(\sigma_{d-1}\), so the standard marking theorem gives the displayed
point-process convergence.
\end{proof}

Put \(\rho_n=\sqrt{2b_{n,d}}\) and
\[
    \lambda_n=\kappa_n\rho_n^2 .
\]
In the critical window \(\lambda_n=O(\log\log n)\), and it is useful to write
\[
    s_n=\lambda_n-(k-1)\log\log n+\log\Gamma(k).
\]
Completing the square in one summand of \(D_n\), or equivalently using
Lemma~\ref{lem:bump-representation}, shows that an observation with
\(Y_i=b_{n,d}+x\) has optimized log contribution \(x-s_n+o_{\Pbb}(1)\),
uniformly for bounded \(x\). Thus an observation is locally supercritical
precisely when its radial excess \(x\) exceeds \(s_n\); these observations
index the candidate KKT crossings.

To resolve each candidate spatially, write the score and observation in normal
coordinates about a direction \(v\). Let
\(\exp_v:T_vS^{d-1}\to S^{d-1}\) be the Riemannian exponential map and put
\[
    t=\left(\rho_n+\frac q{\rho_n}\right)v,\qquad
    z=\left(\rho_n+\frac x{\rho_n}\right)\exp_v\left(\frac y{\rho_n}\right),
    \qquad y\in T_vS^{d-1}.
\]
Then, uniformly for \(x,q,|y|=O(\log\log n)\),
\[
    t\cdot z-\frac{|t|^2}{2}-\log n-\kappa_n|t|^2
    =
    (k-1)\log\log n-\log\Gamma(k)-\lambda_n
    +x-\frac{|y|^2}{2}
    +
    O\!\left(\frac{(\log\log n)^4}{\rho_n^2}\right).
\]
The radial score offset \(q\) cancels to first order. A radial excess \(x\)
therefore raises the optimized contribution by \(x\), whereas an angular
displacement \(y/\rho_n\) lowers it by \(|y|^2/2\); the variance defect
subtracts \(\lambda_n\). The estimates below show that the rest of the sample
contributes \(o_{\Pbb}(1)\) uniformly on each candidate patch.
Subsection~\ref{sec:higher-dimensional-support-counts} uses these patch
estimates to recover the fitted support.

For the uniform exclusion argument, use the centered radial maximum
\[
    M_n^c
    =
    \max_i\left(\frac{|Z_i-\bar Z_n|^2}{2}-b_{n,d}\right).
\]
It differs from \(M_n=\max_i(Y_i-b_{n,d})\) by \(o_{\Pbb}(1)\), since
\[
    |M_n^c-M_n|
    \le
    \max_i|Z_i|\,|\bar Z_n|+\frac{|\bar Z_n|^2}{2}
    =
    o_{\Pbb}(1).
\]

\subsection{Uniform moving-shell estimates}
\label{sec:sharp-high-dimensional-estimates}

This subsection supplies the estimates used to prune subthreshold radial
patches in Lemma~\ref{lem:support-dge3-deleted-field}. Bennett's inequality
controls the center of each angular patch, while scaled derivative bounds and
Morrey's inequality control oscillation across patches of diameter
\(1/\rho_n\). The following bulk estimates handle smaller score radii by a
clipped empirical-process argument.

For the moving-shell bound, replace the field by an independent truncated
field. On the radial-maximum event used in the comparison, every observation
lies below the deterministic truncation radius, so the two fields agree.

Fix a gap \(a>0\).  Set \(r=|t|\),
\[
    T_n(a)=\log n+\lambda_n-a,
    \qquad
    R_n(a)=\{2T_n(a)\}^{1/2},
\]
and write \(R_a=R_n(a)\).  Define
\[
    V_{i,t}^{(a)}
    =
    \frac1n
    e^{-\kappa_nr^2}
    \exp\left\{t\cdot Z_i-\frac{r^2}{2}\right\}
    \mathbf 1_{\{|Z_i|\le R_a\}},
    \qquad
    S_t^{(a)}=\sum_iV_{i,t}^{(a)}.
\]
Let
\[
    \mathcal T_n(B)=
    \left\{t:\ \left||t|^2/2-b_{n,d}\right|\le B\log\log n\right\}.
\]
The centered field is compared to this independent truncated field with a
small deterministic loss in the gap.  Put
\[
    \alpha_n=(\log\log n)^{-2},
    \qquad
    G_n=\{\rho_n|\bar Z_n|\le\alpha_n\}.
\]
Since \(\alpha_n\sqrt n/\rho_n\to\infty\), \(\Pbb(G_n^c)\to0\), and on
\(G_n\), for all large \(n\),
\[
    \sup_{t\in\mathcal T_n(B)}|t\cdot\bar Z_n|\le2\alpha_n .
\]
If \(M_n^c\le s_n-\eta\), then
\(|Z_i-\bar Z_n|^2/2\le\log n+\lambda_n-\eta\) for all \(i\).  On \(G_n\)
this implies \(|Z_i|^2/2\le\log n+\lambda_n-\eta/2\) for all large \(n\).
Consequently, uniformly over \(t\in\mathcal T_n(B)\),
\[
    D_n(t)\le e^{2\alpha_n}S_t^{(\eta/2)} .
\]
It is therefore enough to prove that, for fixed \(a>0\), the independent
truncated field stays below \(1-o(1)\) uniformly on \(\mathcal T_n(B)\).

The first input bounds the moments of one truncated summand. Exponential
tilting turns each moment into the probability that a noncentral Gaussian lies
in the truncation ball; a one-dimensional Gaussian tail estimate then gives
the required decay.
\begin{proposition}
\label{prop:noncentral-ball}
Fix \(d\ge3\), \(C_0<\infty\), \(B<\infty\), \(p>1\), and \(a>0\).  Suppose
the sharp-scale relation
\[
    s_n
    =
    \lambda_n-\left(\frac d2-1\right)\log\log n+\log\Gamma(d/2)
\]
satisfies \(|s_n|\le C_0\).  Uniformly over deterministic sequences satisfying
this relation, over directions \(t/|t|\), and over
\[
    \left|\frac{r^2}{2}-b_{n,d}\right|\le B\log\log n,
\]
one has
\[
    \sum_{i=1}^n\E \{V_{i,t}^{(a)}\}^p
    \le
    C_{p,a,B,C_0,d}\rho_n^{-1}
    \exp\{-\lambda_n-(p-1)a+o(1)\}.
\]
\end{proposition}

\begin{proof}
The exact identity is
\[
    \sum_{i=1}^n\E \{V_{i,t}^{(a)}\}^p
    =
    n^{1-p}
    e^{-p\kappa_nr^2+(p^2-p)r^2/2}
    \Pbb\{|G+pt|\le R_n(a)\},
    \]
where \(G\sim\cN(0,I_d)\).  Align \(t=re_1\), write \(G=We_1+U\), and put
\(\Delta=pr-R_n(a)\).  In the moving window,
\(\Delta=(p-1)\rho_n+O((\log\log n)/\rho_n)\), so
\(\Delta\ge c_p\rho_n\).  If \(|G+pt|\le R_n(a)\), then
\[
    W\le -\left(pr-\sqrt{R_n(a)^2-|U|^2}\right).
\]
Mills' bound and the inequality
\[
    pr-\sqrt{R_n(a)^2-|U|^2}
    \ge
    \Delta+\frac{|U|^2}{2R_n(a)}
\]
give
\[
    \Pbb\{|G+pt|\le R_n(a)\}
    \le
    C_{p,d}\rho_n^{-1}
    \exp\left\{-\frac{(pr-R_n(a))^2}{2}\right\}.
\]
Substitution gives
\[
    \sum_{i=1}^n\E \{V_{i,t}^{(a)}\}^p
    \le
    C_{p,d}\rho_n^{-1}
    \exp\left\{
        -p\kappa_nr^2
        +(p-1)(\lambda_n-a)
        -\frac p2(r-R_n(a))^2
    \right\}.
\]
Since the sharp-scale assumption gives \(\lambda_n=O(\log\log n)\), one has
\(\kappa_nr^2=\lambda_n+o(1)\) uniformly in the displayed shell. This proves
the stated estimate.
\end{proof}

Under the same sharp-scale assumptions, the truncation also gives the
pointwise jump bound
\[
    V_{i,t}^{(a)}\le e^{-a+o(1)}
\]
uniformly in the moving shell.  Indeed
\[
    \log V_{i,t}^{(a)}
    \le
    -\log n-\kappa_nr^2-\frac{r^2}{2}+rR_n(a)
    =
    \lambda_n-a-\kappa_nr^2-\frac12(r-R_n(a))^2.
\]
Thus \(V_{i,t}^{(a)}\le c_a<1\) for large \(n\), where \(c_a\) is deterministic.

The moment and jump bounds give a fixed-location tail strong enough for the
angular union bound after the Morrey step.
\begin{proposition}
\label{prop:fixed-t-bennett}
Fix \(d\ge3\), \(C_0<\infty\), \(B<\infty\), and \(a>0\), and suppose the
sharp-scale relation in Proposition~\ref{prop:noncentral-ball} satisfies
\(|s_n|\le C_0\).  Uniformly over deterministic sequences satisfying this
relation and over the moving shell, for every deterministic
\(\delta_n\to0\) with \(0\le\delta_n\le1/2\) eventually,
\[
    \Pbb\{S_t^{(a)}\ge1-\delta_n\}
    \le
    C_{a,B,C_0,d}
    \left(\rho_n^{-1}e^{-\lambda_n}\right)^{
        e^a(1-\delta_n)-o(1)} .
\]
In particular, at level \(1\) the exponent is \(e^a-o(1)\).
\end{proposition}

\begin{proof}
Let \(\mu_t=\E S_t^{(a)}\), \(v_t=\sum_i\operatorname{Var}(V_{i,t}^{(a)})\),
\(A_n=\rho_n^{-1}e^{-\lambda_n}\), and \(x_n=1-\delta_n\).
Proposition~\ref{prop:noncentral-ball} with \(p=2\) gives
\[
    v_t\le C\rho_n^{-1}e^{-\lambda_n-a+o(1)},
    \qquad
    \mu_t\le e^{-\lambda_n+o(1)}=o(1)
\]
under the displayed \(d\ge3\) sharp-scale assumption.  Bennett's inequality for
\(0\le V_{i,t}^{(a)}\le c_a<1\) gives
\[
    \Pbb\{S_t^{(a)}\ge x_n\}
    \le
    \exp\left[
        -\frac{v_t}{c_a^2}
        h\left(\frac{c_a(x_n-\mu_t)}{v_t}\right)
    \right],
    \qquad h(u)=(1+u)\log(1+u)-u .
\]
Using the sharp jump bound \(V_{i,t}^{(a)}\le e^{-a+o(1)}\) and
\(\log A_n^{-1}=(d-1)\log\rho_n+O(1)\), the Bennett rate is at least
\[
    \{e^a(1-\delta_n)-o(1)\}\log A_n^{-1}.
\]
This proves the displayed estimate. Only the second moment is needed at a
fixed location; the derivative estimates used for the Morrey transfer require
higher moments.
\end{proof}

At the sharp scale \(s_n=O(1)\), equivalently
\(\lambda_n=(d-2)\log\rho_n+O(1)\), the quantity
\(\rho_n^{-1}e^{-\lambda_n}\) is \(\rho_n^{-(d-1)+o(1)}\).  The strict
near-unit exponent \(e^a(1-o(1))>1\) is therefore exactly what is needed to
union bound over \(\rho_n^{d-1+o(1)}\) angular locations.
Throughout the remaining \(d\ge3\) estimates, ``at the sharp scale
\(s_n=O(1)\)'' means uniformly over deterministic sequences satisfying
\(|s_n|\le C\), for each fixed \(C<\infty\); constants may depend on \(C\).

Scaled derivatives control oscillation on angular boxes at scale
\(1/\rho_n\), reducing the continuum to a net of size
\(\rho_n^{d-1+o(1)}\).
\begin{lemma}
\label{lem:scaled-derivative-morrey}
Fix \(d\ge3\), \(p>d\), \(a>0\), and suppose \(s_n=O(1)\).
Let \(Q\subset\R^{d-1}\) be a fixed cube, and let
\(\{u_\alpha\}\) be a bounded-overlap family of smooth charts such that
\(\{u_\alpha(Q/\rho_n)\}\) covers \(S^{d-1}\), the number of charts is
\(O(\rho_n^{d-1})\), and all first two derivatives of the charts are
bounded uniformly in \(\alpha\) and \(n\).  For each fixed \(B<\infty\), in
moving-shell coordinates
\[
    t=t(q,y)=\sqrt{2(b_{n,d}+q)}\,u_\alpha(y/\rho_n),
    \qquad |q|\le B\log\log n,\quad y\in Q,
\]
one has
\[
    \E\sum_\alpha\int_{|q|\le B\log\log n}\int_Q
    |\nabla_{q,y}S_{t(q,y)}^{(a)}|^p\,dy\,dq
    \le(\log\log n)^C .
\]
The displayed derivative-moment bound also holds on every fixed physical-offset annulus
\[
    t=t(h,y)=(R_a-h)u_\alpha(y/\rho_n),
    \qquad |h|\le A,\quad y\in Q,
\]
with \(C=C(a,A,B,d,p)\). Consequently, in either region and for every fixed
\(\gamma>0\), a net of cardinality
\[
    \rho_n^{d-1}(\log\log n)^{C_\gamma}=\rho_n^{d-1+o(1)}
\]
suffices, with probability tending to one, to approximate \(S_t^{(a)}\)
up to the prescribed error
\[
    \varepsilon_{n,\gamma}=(\log\log n)^{-\gamma}.
\]
\end{lemma}

\begin{proof}
The truncation set \(\{|Z_i|\le R_a\}\) is independent of \(t\), so
\(S_t^{(a)}\) is \(C^1\) in the displayed coordinates.  For any coordinate
\(\xi\in\{q,y_1,\ldots,y_{d-1}\}\) or
\(\xi\in\{h,y_1,\ldots,y_{d-1}\}\),
\[
    \partial_\xi V_{i,t}^{(a)}
    =
    V_{i,t}^{(a)}M_\xi(Z_i,t),
    \qquad
    M_\xi(z,t)=\{z-(1+2\kappa_n)t\}\cdot\partial_\xi t .
\]
The coordinate choices give
\[
    |\partial_qt|=O(\rho_n^{-1}),\qquad
    |\partial_ht|=1,\qquad
    |\partial_{y_j}t|=O(1),\qquad
    t\cdot\partial_{y_j}t=0 .
\]
For \(m\in\{2,p\}\), exponential tilting gives the exact identity
\[
    \sum_i\E|\partial_\xi V_{i,t}^{(a)}|^m
    =
    n^{1-m}e^{-m\kappa_nr^2+(m^2-m)r^2/2}
    \E\!\left[
        |M_\xi(G+mt,t)|^m\mathbf 1_{\{|G+mt|\le R_a\}}
    \right].
\]
We next bound the tilted expectation.  Align \(t=re_1\), write
\(z=G+mt\), and decompose
\[
    z=z_1e_1+z_\perp,\qquad
    \zeta=R_a-z_1,\qquad \Delta_m=mr-R_a .
\]
In both the moving shell and fixed physical-offset annuli,
\(\Delta_m=(m-1)\rho_n+O(\log\log n)\), and hence
\(\Delta_m\ge c_m\rho_n\) for large \(n\).  We claim that, for every fixed
integer \(L\),
\[
    \E\!\left[
        (1+\zeta_+ + |z_\perp|)^L\mathbf 1_{\{|z|\le R_a\}}
    \right]
    \le
    C_{L,m}\rho_n^{-1}e^{-\Delta_m^2/2}.
\]
Indeed, write \(G=We_1+U\).  If \(|G+mt|\le R_a\), then
\[
    W\le -\Delta_m-\frac{|U|^2}{2R_a},
\]
because \(\sqrt{R_a^2-|U|^2}\le R_a-|U|^2/(2R_a)\).  The one-dimensional
Gaussian tail-moment bound
\[
    \E\{(1+(-x-W)_+)^L\mathbf 1_{W\le -x}\}
    \le C_Lx^{-1}e^{-x^2/2},\qquad x\ge1,
\]
with \(x=\Delta_m+|U|^2/(2R_a)\), gives the claim after integration in
\(U\): since \(\Delta_m/R_a\) stays bounded away from zero and infinity,
the factor \(e^{-x^2/2}\) is at most
\(e^{-\Delta_m^2/2}e^{-c|U|^2}\).

In these coordinates the derivative marks have the following deterministic
bounds on \(\{|z|\le R_a\}\):
\[
\begin{aligned}
    M_q(z,t)
    &=
    \{z-(1+2\kappa_n)t\}\cdot\partial_qt
      =O\{(1+\zeta_+)/\rho_n+\kappa_n\rho_n\},\\
    M_h(z,t)
    &=
    \{z-(1+2\kappa_n)t\}\cdot\partial_ht
      =O_A(1+\zeta_++\kappa_n\rho_n),\\
    M_{y_j}(z,t)
    &=
    \{z-(1+2\kappa_n)t\}\cdot\partial_{y_j}t
      =O\{1+|z_\perp|+(\log\log n)(1+\zeta_+)/\rho_n\}.
\end{aligned}
\]
For \(M_q\), use \(R_a-r=O(\log\log n/\rho_n)\) in the moving shell; for
\(M_h\), use \(|h|\le A\); for \(M_{y_j}\), use
\(t\cdot\partial_{y_j}t=0\) and the uniform chart derivative bounds.  Since
\(\kappa_n\rho_n=O(\log\log n/\rho_n)=o(1)\), the tilted tail-moment estimate
from Proposition~\ref{prop:noncentral-ball} shows that the derivative marks cost only fixed powers
of \(\log\log n\).  Combining this with the same algebra as in
Proposition~\ref{prop:noncentral-ball} yields, uniformly in the moving shell
and in fixed physical-offset annuli,
\[
    \sum_i\E|\partial_\xi V_{i,t}^{(a)}|^m
    \le
    (\log\log n)^C\rho_n^{-1}e^{-\lambda_n}
    \le
    (\log\log n)^C\rho_n^{-(d-1)} .
\]
Here we used \(e^{-\lambda_n}\le C\rho_n^{-(d-2)}\), which follows from
\(s_n=O(1)\).

Apply Rosenthal's inequality to
\(\partial_\xi S_t^{(a)}=\sum_i\partial_\xi V_{i,t}^{(a)}\).  The centered
\(p\)th-moment term is bounded by the display with \(m=p\), and the variance
term is bounded by the display with \(m=2\) raised to the power \(p/2\),
which is smaller.  It remains only to bound the mean.  Rotational symmetry
makes the angular means zero.  For radial derivatives, write
\[
    \E S_t^{(a)}=e^{-\kappa_nr^2}P_r,\qquad
    P_r=\Pbb\{|G+re_1|\le R_a\}.
\]
The derivative \(P_r'\) is uniformly bounded, because
\[
    P_r'=\int_{|z|\le R_a}(z_1-r)\phi_d(z-re_1)\,dz
\]
and the absolute value of this integral is at most \(\E|G_1|\).  Hence
\(\partial_q\E S_t^{(a)}=O(e^{-\lambda_n}/\rho_n)\), since
\(\partial_q r=1/r=O(\rho_n^{-1})\) and
\(\kappa_n=O(\log\log n/\rho_n^2)\).  On a fixed physical-offset annulus,
\(\partial_h\E S_t^{(a)}=O_A(e^{-\lambda_n})\).  Their \(p\)th powers are
also \(O(\rho_n^{-(d-1)})\), because \(p>d\) and \(d\ge3\).  Therefore
\[
    \E|\partial_\xi S_t^{(a)}|^p
    \le(\log\log n)^C\rho_n^{-(d-1)}
\]
uniformly in all coordinates and all displayed parameter values.

Integrating this pointwise bound over \(O(\rho_n^{d-1})\) angular charts and
over \(q\)-length \(O(\log\log n)\), or fixed \(h\)-length \(O_A(1)\), gives
the displayed integral bound.  By Markov, the total gradient integral is at
most \((\log\log n)^{C'}\) with probability tending to one.  Fix
\(\gamma>0\), choose
\[
    K>\frac{\gamma+C'/p}{1-d/p},
\]
and, on this event, partition every coordinate patch into \(d\)-dimensional
cubes.  Take the side length to be
\[
    \ell_n=(\log\log n)^{-K}.
\]
The Morrey--Poincare inequality on each cube gives
\[
    O\!\left(\ell_n^{1-d/p}(\log\log n)^{C'/p}\right)
    =
    O\!\left((\log\log n)^{-K(1-d/p)+C'/p}\right),
\]
as a uniform bound for the oscillation of \(S_t^{(a)}\) on that cube.  For
the displayed choice of \(K\), this is at most
\((\log\log n)^{-\gamma}\).  Taking one point from each cube gives a net of
size \(\rho_n^{d-1}(\log\log n)^C\) in either coordinate region.
\end{proof}

The Morrey estimate is used below only to transfer fixed-location bounds to
the angular continuum. Its net has size \(\rho_n^{d-1+o(1)}\), which is small
enough for the Bennett exponents in Proposition~\ref{prop:fixed-t-bennett}.

For radii below the moving shell, use physical distance from the truncated
sample edge. Put
\[
    R_a=R_n(a),\qquad h=R_a-|t|.
\]
A bounded interval of \(h\) corresponds to a radial-energy interval of length
\(O(\rho_n)\), so a direct energy-coordinate net would be unnecessarily large.
The fixed-location estimates retain the Gaussian penalty in \(h\).
\begin{lemma}
\label{lem:physical-offset-fixed-t}
Fix \(d\ge3\), \(a>0\), and \(A<\infty\), and suppose \(s_n=O(1)\).
Uniformly over directions and \(|h|\le A\), for each fixed \(p>1\),
\[
 V_{i,t}^{(a)}\le e^{-a-h^2/2+o(1)},\qquad
 \sum_i\E\{V_{i,t}^{(a)}\}^p
 \le C_{p,a,A,d}\rho_n^{-1}
 e^{-\lambda_n-(p-1)a-ph^2/2+o(1)}.
\]
Consequently, for every fixed \(\zeta\in(0,1)\),
\[
 \Pbb\{S_t^{(a)}\ge\zeta\}
 \le \left(\rho_n^{-1}e^{-\lambda_n}\right)^{
 e^{a+h^2/2}\zeta-o(1)}.
\]
\end{lemma}

\begin{proof}
The completing-square and Mills-ratio calculation used in
Proposition~\ref{prop:noncentral-ball} sharpens in the physical coordinate
\(h=R_a-r\).  Uniformly for bounded \(h\) and every fixed \(p>1\),
\[
    V_{i,t}^{(a)}\le \exp\{-a-h^2/2+o(1)\},
    \qquad
    \sum_i\E\{V_{i,t}^{(a)}\}^p
    \le
    C\rho_n^{-1}\exp\{-\lambda_n-(p-1)a-ph^2/2+o(1)\}.
\]
The first bound is the completing-square inequality
\[
    \log V_{i,t}^{(a)}
    \le
    \lambda_n-a-\kappa_nr^2-\frac12(R_a-r)^2
    =
    -a-h^2/2+o(1).
\]
For the second bound, exponential tilting reduces
\(\sum_i\E\{V_{i,t}^{(a)}\}^p\) to a Gaussian probability of the ball
\(\{|G|\le R_a\}\) under a mean shifted by \(pr\) in the \(t\)-direction.
The nearest boundary distance is
\[
    pr-R_a=(p-1)R_a-ph.
\]
Using the Mills bound as in Proposition~\ref{prop:noncentral-ball} gives
\[
    \Pbb\{|G+pt|\le R_a\}
    \le
    C_{p,A,d}\rho_n^{-1}
    \exp\left\{-\frac{(pr-R_a)^2}{2}\right\}.
\]
Substituting \(r=R_a-h\), \(R_a^2/2=\log n+\lambda_n-a\), and
\(\kappa_nr^2=\lambda_n+o(1)\) into the exact tilted identity gives
\[
\begin{aligned}
&n^{1-p}e^{-p\kappa_nr^2+(p^2-p)r^2/2}
    \exp\left\{-\frac{(pr-R_a)^2}{2}\right\}  \\
&\qquad\le
    \exp\{-\lambda_n-(p-1)a-ph^2/2+o(1)\},
\end{aligned}
\]
which is the displayed moment estimate for \(p>1\).  Separately,
\[
    \mu_t=\E S_t^{(a)}
    =
    e^{-\kappa_nr^2}\Pbb\{|G+t|\le R_a\}
    \le e^{-\kappa_nr^2}
    =
    e^{-\lambda_n+o(1)}
    =
    o(1).
\]
The \(p=2\) case gives
\[
    v_t:=\sum_i\operatorname{Var}(V_{i,t}^{(a)})
    \le C\rho_n^{-1}e^{-\lambda_n-a-h^2+o(1)},
\]
and the jump bound is \(V_{i,t}^{(a)}\le e^{-a-h^2/2+o(1)}\). Bennett's
inequality gives
\[
    \Pbb\{S_t^{(a)}\ge\zeta\}
    \le
    \left(\rho_n^{-1}e^{-\lambda_n}\right)^{e^{a+h^2/2}\zeta-o(1)}.
\]
\end{proof}

For growing physical offset \(h=R_a-|t|\), the fixed-location tail becomes
much smaller.
\begin{lemma}
\label{lem:growing-offset-fixed-t}
Fix \(d\ge3\), \(a>0\), \(A_0>0\), and \(\tau\in(0,1)\).  At the sharp
scale \(s_n=O(1)\), uniformly over directions and over
\[
    h=R_a-|t|\in[A_0,\rho_n/2+2],
\]
one has
\[
    \mu_t:=\E S_t^{(a)}=o(1),\qquad
    \max_i V_{i,t}^{(a)}\le \beta_h:=\exp\{-a-h^2/4\},
\]
and
\[
    v_t:=\sum_i\operatorname{Var}(V_{i,t}^{(a)})
    \le
    \beta_h\,\rho_n^{-(d-2)+o(1)} .
\]
Consequently,
\[
    \Pbb\{S_t^{(a)}\ge\tau\}
    \le
    \exp\{-c_{\tau,d}e^{a+h^2/4}\log\rho_n\}
\]
for all large \(n\), with \(c_{\tau,d}>0\).
\end{lemma}

\begin{proof}
Write \(R=R_a\), \(r=|t|=R-h\), and use \(R=\rho_n+O((\log\log n)/\rho_n)\),
\(\lambda_n=\kappa_n\rho_n^2=O(\log\log n)\).  Since \(h\ge A_0\),
\[
    \kappa_n(\rho_n^2-r^2)
    \le
    C\lambda_n h/\rho_n
    \le h^2/4
\]
uniformly on the displayed range, for all large \(n\).  The pointwise
completing-square bound therefore gives
\[
    \log V_{i,t}^{(a)}
    \le
    \lambda_n-a-\kappa_nr^2-h^2/2
    =
    -a-h^2/2+\kappa_n(\rho_n^2-r^2)
    \le -a-h^2/4 .
\]
This proves the jump bound.

The mean satisfies
\[
    \mu_t
    =
    e^{-\kappa_nr^2}\Pbb\{|G+t|\le R\}
    \le
    e^{-\kappa_nr^2}
    \le
    e^{-\lambda_n/10}=o(1),
\]
because \(r\ge R-\rho_n/2-2\ge2\rho_n/5\) for large \(n\), and
\(\lambda_n=(d-2)\log\rho_n+O(1)\).

It remains to bound the variance.  The exact second-moment identity is
\[
    \sum_i\E\{V_{i,t}^{(a)}\}^2
    =
    n^{-1}e^{-2\kappa_nr^2+r^2}
    \Pbb\{|G+2t|\le R\}.
\]
If \(h\le R/2-1\), then the ball event implies a one-dimensional
left-tail event of size \(R-2h\).  Hence
\[
    \Pbb\{|G+2t|\le R\}
    \le
    C\exp\{-(R-2h)^2/2\},
\]
and the algebra
\[
    -\log n+r^2-(R-2h)^2/2
    =
    \lambda_n-a-h^2
\]
gives
\[
    \sum_i\E\{V_{i,t}^{(a)}\}^2
    \le
    C\exp\{-a-h^2-\lambda_n
        +2\kappa_n(\rho_n^2-r^2)\}
    \le
    C e^{-a-h^2/2}\rho_n^{-(d-2)+o(1)} .
\]
This is bounded by \(\beta_h\rho_n^{-(d-2)+o(1)}\).  If
\(h>R/2-1\), then \(\Pbb\{|G+2t|\le R\}\le1\) and
\[
    n^{-1}e^{-2\kappa_nr^2+r^2}
    \le
    \exp\{-(1/2+o(1))\log n\},
\]
because \(r\le R/2+1\).  Since \(h\le\rho_n/2+2\), this is again
\(\le\beta_h\rho_n^{-(d-2)+o(1)}\): here the left side is
\(\exp\{-(1/2+o(1))\log n\}\), whereas \(\beta_h^{-1}\) is at most
\(\exp\{(1/8+o(1))\log n\}\).  The variance bound follows.

For large \(n\), \(\mu_t\le\tau/2\).  Bennett's inequality for centered
summands bounded by \(\beta_h\) and total variance \(v_t\) yields
\[
    \Pbb\{S_t^{(a)}\ge\tau\}
    \le
    \exp\left[
        -\frac{v_t}{\beta_h^2}
        H\!\left(\frac{\beta_h(\tau-\mu_t)}{v_t}\right)
    \right],
    \qquad
    H(u)=(1+u)\log(1+u)-u .
\]
Since \(\beta_h(\tau-\mu_t)/v_t\ge c_\tau\rho_n^{d-2-o(1)}\), the
exponent is at least \(c_{\tau,d}\beta_h^{-1}\log\rho_n\).  This is the
claimed tail bound.
\end{proof}

The exponential decay in \(h^2\) permits a union bound over a polynomial-size
Euclidean net. The next lemmas handle the remaining bulk
\(|t|\le\rho_n/2\).
